\section{State of the Art}

There are already hundreds of smoking cessation apps available, but unfortunately, most apps are not scientifically validated \cite{boldSmartphoneAppsSmoking2023}. In a study by Bold et al., only 11 apps had published results supporting their efficacy.




\subsection{Feasibility of apps}

% \textbf{How feasible is it to use apps? What makes an app effective? Maybe go into detail about how behavior change in general is hard (see Dries: only about 30\% success rate, even for tobaccologists).}

While many apps are available, the effectiveness of these apps may be limited \cite{mershaEffectivenessTechnologybasedInterventions2021,abromsContentAnalysisPopular2013}. Luckily, it has been shown that digital interventions, through apps, can work \cite{mershaEffectivenessTechnologybasedInterventions2021, PersonalizedDigitalTools2025a}, when the right interventions are used.
Even people who already make use of face-to-face care can find benefits in using digital help \cite{mershaEffectivenessTechnologybasedInterventions2021}.

\subsection{Description of existing apps}

% \textbf{For some of the most salient apps: describe the methods they use. Are the methods scientifically based? Do they rely on scientific evidence?}

From the literature, we see that a lot of apps do not use science-based interventions \cite{goldgofEfficacyQuitSureApp2024}.

To collect information about the most popular smoking cessation apps, we searched the internet for the top-rated apps, and looked at Facebook groups to see what users recommend to each other.

We looked at a few of the most popular smoking cessation apps, and we saw that they often employ these methods:
\begin{itemize}
    \item \textbf{Progress trackers.} There are many variations of this; some track the number of days without smoking, some track the amount of money saved by not smoking, some track lifetime saved by not smoking.
    \item \textbf{Social aspects.} Some apps use a social aspect, like forums or other chat functions.
    \item \textbf{Planners.} Apps might help users plan designated smoke-free days, or plan a quit date.
    \item \textbf{\ac{cbt}}. Some apps use \ac{cbt} as a basis for their interventions.
    \item \textbf{Gamification.} Many apps use badges, mileposts, or levels to signify progress. Some apps also include games as a distraction for when a craving hits.
\end{itemize}

To make things more concrete, we will now look at a subset of the most popular apps, which we believe are scientifically based.

\subsubsection{QuitSure}

This app has published a study, showing its effectiveness \cite{BestQuitSmoking, goldgofEfficacyQuitSureApp2024}. The study reports that 80.1\% of the users maintained abstinence for at least 30 days after program completion \cite{goldgofEfficacyQuitSureApp2024}.

The app is based on techniques from psychology, \ac{cbt}, and mindfulness \cite{goldgofEfficacyQuitSureApp2024}.

In an onboarding process, the user is asked to think about their smoking quantity, patterns, past quitting attempts, and relapses. The app then assesses if the user is ready to quit. After onboarding, the user is intended to follow a program consisting of content and video exercises. This program requires 6 to 10 hours over 6 to 12 days to complete \cite{goldgofEfficacyQuitSureApp2024}.

In their discussion, the researchers note that there is a high dropout rate of 59.4\% \cite{goldgofEfficacyQuitSureApp2024}, as a result of the high engagement requirement. 
Additionally, most relapses were a result of high cravings, for which the app did not have specific interventions after program completion.

Unfortunately, the study does not go into detail on what specific modules in the app contributed most to successful smoking cessation.




\subsection{Behavior models}

% \textbf{Found: HAPA, COM-B, TDF. Explain what each model does, on what timescale, and investigate how these models can be merged / used together.}

Through our first meeting with Geert Crombez, we already came into contant with the \ac{hapa} and \ac{bct} models.
In our own search, we came across the \ac{comb} and \ac{tdf} models.

These models all serve a slightly different purpose and operate on different timescales.

The \ac{comb} model [citation needed] describes the relationship between capability, opportunity, and motivation in influencing behavior. This means you can influence behavior by tweaking the other factors.
The model can thus describe what factors you need to change.

\ac{hapa} \cite{schwarzerModelingHealthBehavior2008,schwarzerHealthActionProcess2016} makes a distinction between a preintentional phase and a volitional phase. In the first phase, a person develops their motivation and plan of action. Once this is "good intention" is sufficiently developed, the desired health action can be initiated.
Now, in the volitional phase, this action needs to be maintained, which needs self-regulatory skills and strategies \cite{schwarzerModelingHealthBehavior2008}.

Compared to \ac{comb}, \ac{hapa} focuses more on the psychological processes that drive behavior change, particularly the transition from intention to action and the maintenance of that action, rather than the physical or environmental factors.

We want to use a hybrid approach, combining models to find the best intervention for the user for each stage in the program. Further research is required into how the models interact, and determine the feasibility of this idea.




\subsection{Related research topics}

% \textbf{Related topics: how to use chatbot integration correctly? How to steer our interventions according to user feedback (either direct or indirect feedback)?}

In addition to the models discussed, we also need to consider the integration of chatbot technology and the use of user feedback to improve our interventions.

To successfully achieve this, we need to perform more research into the effectiveness of chatbots, when to use them, and when not to.

Additionally, we want to look into how we can effectively learn to steer our interventions, when the user provides feedback. We want to explore at least two paths: how to learn from indirect feedback, like swiping away or ignoring a notification, and how to learn from direct feedback, like a yes-or-no question we ask of the user.
